\chapter{Introduction}
\label{ch:introduction}

Climate change is increasing both the frequency and severity of wildfires by altering temperature, precipitation, and vegetation patterns \cite{dennison2009climate}, and burnt area in EU countries has repeatedly exceeded historical averages in recent fire seasons \cite{effis2024}. Aerial firefighting (AFF) remains one of the primary tools available to suppress wildfires before they escalate, yet the coordination of aircraft fleets during an active fire is still largely dependent on the experience and judgement of human operators. Given the scale, time pressure, and safety stakes involved, there is a clear need for decision-support tools that can compute efficient aircraft routing and suppression strategies fast enough to be usable in live operations.

The Vehicle Routing Problem (VRP) and its variants have long been adapted to disaster and humanitarian relief settings, where capacity limits, multiple depots, and split deliveries are common features of the underlying logistics, as in the Split Delivery VRP \cite{dror1989split} and drone-based relief distribution under payload and energy constraints \cite{rabta2018drone}. \citet{chalabi2024affvrp} adapt these ideas specifically to wildfire suppression, introducing the Aerial FireFighting Vehicle Routing Problem (AFFVRP): a capacitated, multi-trip VRP with time windows and a hierarchical objective that first minimises time-to-last-drop and then total flight time, using a heterogeneous fleet of tankers and scoopers. \citet{liu2025hierarchical} extend this direction with a hierarchical decision-making model for multi-aircraft \emph{formations}, explicitly modelling the water-release process and tactical coordination behaviours that pure routing formulations do not capture. These approaches motivate the joint consideration of aircraft logistics and the suppression actions through which water is delivered to a fire.

Physically grounded fire spread models address the evolution of the fire itself. FARSITE \cite{finney1998farsite} propagates a fire perimeter using Huygens' principle of wavelet expansion driven by the Rothermel surface fire spread equations. DEVS-FIRE \cite{hu2012devsfire} instead represents the burning area through discrete-event cells that can respond to suppression inputs. Connecting a routing model to either representation requires more than transferring a list of target locations: aircraft actions must produce explicit changes in the fire state, and those changes must be accounted for consistently when the next set of decisions is made.

This dissertation develops an optimisation framework for aerial firefighting that discretises the operational area into a finite grid of cells and a finite set of timesteps. Each cell carries one of three states --- unburned, burning, or extinguished --- and aircraft routing, water capacity, and suppression decisions are expressed over the same spatial and temporal representation. The broader framework is intended to support coupling with FARSITE in a rolling-horizon loop. Before that coupling can be evaluated, however, the routing and suppression model must be verified under conditions in which the expected outcome is known independently of the optimiser.

The implementation examined in this dissertation therefore establishes a deterministic \emph{static} experiment. Static refers to a fixed initial fire snapshot: there is no fire propagation, natural burnout, or stochastic change in the fire state during the optimisation horizon. Initially burning cells remain burning until sufficient water has been deposited on them. Aircraft remain heterogeneous, with tankers refilling at an airfield and scoopers at configured water sources, and water drops follow a cell-integrated hexagonal footprint with a Gaussian lateral profile. Suppression is determined by cumulative delivered litres rather than by a probabilistic extinguishment event.

The resulting mixed-integer linear program is solved with Gurobi using two sequential objectives. The first minimises the sum of burning cells across the sampled time boundaries, measured in burning cell-steps. The second minimises total Manhattan flight distance without worsening the primary incumbent. These objectives distinguish the static experiment from the time-to-last-drop and total-flight-time hierarchy of \citet{chalabi2024affvrp}. In particular, reducing burning cell-steps rewards earlier suppression even when the entire fire cannot be extinguished within the available horizon.

Verification combines independent schedule replay with five controlled cases whose optimal schedules can be derived by hand. These cases isolate a single drop, cumulative suppression across two drops, tanker refilling, scooper refilling, and a fire that cannot be reached under the discrete movement rule. Further regression checks address water conservation, terminal actions, aircraft occupancy, footprint geometry, and geographic input alignment. Small-case optimality is considered separately from the quality of time-limited solutions on the full landscape.

A separate rolling-horizon pipeline provides the wider setting for future integration with FARSITE, run inside FlamMap. The static results reported here do not establish the correctness or operational performance of that dynamic coupling. Further work includes carrying aircraft state across successive windows, evaluating suppression feedback to fire spread, calibrating the suppression threshold against physical evidence, and introducing a discrete-event simulation layer using salabim \cite{vanderham2018salabim}. More detailed aircraft flight dynamics can then be assessed within that layer rather than assumed to follow from the present grid-based movement constraints.

The remainder of the material is organised as follows. Chapter~\ref{ch:literature} reviews fire propagation, water-drop dynamics, and aircraft routing. Chapter~\ref{ch:method} presents the methodology and mathematical formulation of the static experiment. Chapter~\ref{ch:static-validation} describes its verification procedure, reports the recorded validation outcomes, and examines the limitations of the full-landscape runs.
